Large vision-language models (VLMs) often produce fluent answers that are not supported by the image they describe. This thesis studies the fine-grained version of the problem posed by the SHROOM-Visions 2026 shared task: given an image, a question and a model answer in English, French, Italian or Chinese, mark every hallucinated character span, assign it one of five categories (invention, mischaracterization, OCR, miscounting, other), and attach a probability that tracks how many human annotators would mark it.

We propose a detector built from four parts. (i) A Qwen3-VL-8B judge fine-tuned with QLoRA writes the hallucinated spans as JSON. (ii) A label-aware span decoder maps that text back to character offsets while keeping text that annotators gave two labels. (iii) A self-consistency voter combines eight sampled decodes and the greedy decode of the judge by a per-label, per-character vote. (iv) White-box uncertainty probes read the token entropy, likelihood, attention look-back ratio and hidden state of two frozen VLMs from different families, Qwen3-VL-8B and Gemma-3-12B, and decode spans with a BiGRU encoder and five jointly trained binary CRFs. Each probe enters the vote as one more voter. The fused output keeps a graded per-character probability rather than one value per span.

All design choices were made under a pre-registered protocol. Validation was split into a selection half and a confirmation half grouped by image. Each decision compared at most five candidates. Claims needed three training seeds and a paired bootstrap. Averaged over three seeds on the official test set, the detector reaches Cor+Lbl 0.4312, against 0.4132 for the same judge fine-tuned on the original data alone (+0.0180). It improves every seed and every language; IoU rises by 0.021 and Cor by 0.015. On the post-competition leaderboard, its best seed places 5th, 5th, 4th and 2nd in English, French, Italian and Chinese, up from 8th, 6th, 7th and 5th for our shared-task entry.

A second line of experiments tested synthetic training data in isolation. Five conditions were tested, including real-span transplants that a text classifier cannot tell from real spans. None beat the original-data baseline. A clean/labeled decomposition of the metric shows why: synthetic records raise the share of hallucinated training examples, and the gain they bring on hallucinated responses is lost to false alarms on clean ones.
